\documentclass[10pt,conference,compsocconf]{IEEEtran}

\usepackage{hyperref}
\usepackage{graphicx}
\usepackage{booktabs}
\usepackage{amsmath}

\begin{document}
\title{Predicting Coronary Heart Disease from a Health Survey\\with Linear Models}

\author{
  L\'eon Lhuillier, [Teammate 2], [Teammate 3]\\
  \textit{CS-433 Machine Learning, EPFL, Switzerland}
}

\maketitle

\begin{abstract}
We predict whether a person has had a heart attack or coronary heart disease
from 321 answers to the BRFSS 2015 telephone survey, using only linear models
implemented in NumPy. Most of the work is in turning survey \emph{codes} into
meaningful numbers: decoding fields that pack a unit inside the answer,
separating ``refused'' from ``never'', and one-hot encoding categories.
Because only 8.8\% of the people are positive, we optimize the decision
threshold for F1. With an honest 5-fold cross-validation, in which every
statistic and the threshold are fitted on training folds only, regularized
logistic regression reaches an F1 of $0.426\pm0.005$, against $0.415$ for a
naive median-impute-and-standardize baseline.
\end{abstract}

\section{Introduction}

The Behavioral Risk Factor Surveillance System (BRFSS) is a US telephone
survey on health behaviours. We receive 328{,}135 labelled respondents and
109{,}379 test respondents, each described by 321 columns, and must predict a
binary label (heart disease $+1$, otherwise $-1$). Predictions are scored with
the F1 score, $F_1 = 2\,TP/(2\,TP+FP+FN)$. Since only 8.8\% of training
labels are positive, a classifier that always predicts $-1$ is 91.2\%
accurate but has $F_1=0$: accuracy is not a useful guide here.

The project imposes the six linear methods seen in class (least squares,
ridge regression, gradient descent (GD) and stochastic GD on the mean squared
error (MSE), logistic and regularized logistic regression) and forbids external
libraries. Our contribution is therefore mainly a careful, reproducible
preprocessing pipeline and a leakage-free validation protocol.

\section{Data and preprocessing}
\label{sec:prep}

\textbf{Survey codes.} Answers are codes, not measurements. Across most
questions, 7/77/777 mean ``don't know'' and 9/99/999 ``refused'', which we
map to missing (NaN), while 8/88/888 mean ``none'' or ``never'', which we map
to $0$. We wrote a metadata table (\texttt{data\_columns.csv}) giving, for 263
columns, this value mapping, the valid range and whether the column is
categorical. Values outside the valid range become NaN.

\textbf{Special columns.} Twenty columns pack a unit inside the number:
\texttt{FRUIT1}$=205$ means ``5 times per week'', \texttt{EXERHMM1}$=130$
means 1\,h\,30, and \texttt{WEIGHT2} is in pounds unless it starts with 9
(kilograms). We decode them into a single unit per column (times per day,
minutes, kg, cm, \dots) \emph{before} applying the value mapping, so that
codes like 888 are still raw when decoded. A multiple-choice column
(\texttt{PCDMDECN}, ``12'' = options 1 and 2) is split into one binary column
per option.

\textbf{Feature construction.} The following steps are fitted on training data
only (Table~\ref{tab:prep}).
(1)~We drop identifiers and interview metadata (record number, date, state,
phone type, survey weights) and the 58 columns without metadata.
(2)~We drop columns with more than 90\% missing values. Many questions are
only asked to a sub-group (e.g.\ diabetics), so a high missing rate is common:
the median column is 36\% missing.
(3)~Categorical columns with 3 to 15 categories are one-hot encoded with the
categories seen in training; otherwise a value such as ``marital status $=6$''
would be treated as six times status 1.
(4)~For numerical columns, a $\log(1+x)$ transform is applied to non-negative
columns whose 99th percentile exceeds ten times the median (e.g.\ blood-sugar
checks per year, up to 36{,}135).
(5)~Missing values are replaced by the training median, and a binary
``was missing'' indicator is added for every column with more than 1\%
missing values, since a skipped question carries information.
(6)~Columns are standardized with the training mean and standard deviation,
which makes gradient descent well-conditioned and the ridge penalty treat all
features equally. A bias column is added.

\begin{table}[htbp]
  \centering
  \caption{Features after preprocessing (full training set).}
  \label{tab:prep}
  \begin{tabular}{lr}
    \toprule
    Block & Columns \\
    \midrule
    Numerical (median-imputed, standardized) & 139 \\
    \quad of which $\log(1+x)$-transformed & 16 \\
    One-hot (31 categorical variables) & 145 \\
    Missing-value indicators & 109 \\
    Bias & 1 \\
    \midrule
    Total (constant columns removed) & 394 \\
    \bottomrule
  \end{tabular}
\end{table}

\section{Models and validation}

\textbf{Models.} Linear models output a score $s = \tilde{x}^\top w$; we
predict $+1$ when $s\ge t$. For MSE-based models we regress $y\in\{0,1\}$;
logistic models use the same labels. Regularized logistic regression minimizes
the mean negative log-likelihood plus $\lambda\|w\|^2$ by full-batch gradient
descent (step size $\gamma=0.5$, 300 iterations, $w_0=0$). We also tested a
class-weighted variant in which positive samples weigh 3 or 6 times more.

\textbf{Threshold.} With 8.8\% positives, almost no score reaches the
default threshold of 0.5, which gives F1 $\approx 0.13$. We therefore pick $t$
by maximizing F1 over 200 quantiles of the training scores.

\textbf{Cross-validation.} We use 5-fold cross-validation (seed 1). In each
fold, the whole pipeline (all preprocessing statistics, the weights and the
threshold $t$) is fitted on the four training folds, and F1 is measured once
on the held-out fold. Choosing $t$ on the evaluated data, as in our first
experiments, gives optimistic scores. We report the mean and standard
deviation over folds.

\section{Results}

\begin{table}[htbp]
  \centering
  \caption{5-fold cross-validated F1 (mean $\pm$ std). ``Baseline'':
  raw codes, median imputation and standardization only. ``Ours'':
  pipeline of Section~\ref{sec:prep}.}
  \label{tab:results}
  \begin{tabular}{lcc}
    \toprule
    Model & Baseline & Ours \\
    \midrule
    Least squares & $0.397 \pm 0.028$ & singular \\
    Ridge, $\lambda=10^{-4}$ & $0.411 \pm 0.004$ & $0.420 \pm 0.006$ \\
    Ridge, $\lambda=10^{-2}$ & $0.410 \pm 0.004$ & $0.418 \pm 0.006$ \\
    Reg.\ logistic, $\lambda=10^{-4}$ & $0.415 \pm 0.004$ & $\mathbf{0.426 \pm 0.005}$ \\
    Weighted logistic, $\times 3$ & $0.415 \pm 0.004$ & $0.425 \pm 0.005$ \\
    Weighted logistic, $\times 6$ & $0.413 \pm 0.004$ & $0.422 \pm 0.005$ \\
    \bottomrule
  \end{tabular}
\end{table}

Table~\ref{tab:results} shows three findings.
First, preprocessing improves every model by about 0.01 F1, roughly twice the
fold-to-fold standard deviation.
Second, logistic regression beats the MSE-based models by about 0.006. Least
squares is unstable on raw codes (std 0.028) and fails on our features: the
one-hot blocks are linearly dependent, so $X^\top X$ is singular. Ridge
regression removes the problem through the term $\lambda I$.
Third, class weighting does not help once the threshold is tuned: moving $t$
already trades precision for recall, and weighting only shifts the scores.

\begin{table}[htbp]
  \centering
  \caption{Ablation: 5-fold F1 of ridge ($\lambda=10^{-4}$) when one step of
  our pipeline is removed or changed (full pipeline: $0.420\pm0.006$).}
  \label{tab:ablation}
  \begin{tabular}{lcc}
    \toprule
    Variant & F1 & Features \\
    \midrule
    Full pipeline & $0.420 \pm 0.006$ & 394 \\
    No one-hot (categories as numbers) & $0.417 \pm 0.006$ & 301 \\
    Drop columns $>50\%$ missing (not $>90\%$) & $0.417 \pm 0.007$ & 324 \\
    No missing-value indicators & $0.419 \pm 0.007$ & 285 \\
    No $\log(1+x)$ transform & $0.419 \pm 0.007$ & 394 \\
    \bottomrule
  \end{tabular}
\end{table}

The ablation (Table~\ref{tab:ablation}) shows that one-hot encoding and
keeping sparsely answered questions contribute most (about $0.003$ each),
while missing indicators and the $\log$ transform are within noise. No single
step explains the gain on its own: it comes from cleaning the codes as a
whole.

Our submission (\texttt{run.py}) refits the pipeline and the regularized
logistic regression ($\lambda=10^{-4}$) on the full training set, with the
threshold chosen on the training scores.

\section{Discussion}

The decisive steps were cheap: tuning the threshold for F1 (from 0.13 to about
0.42), then making the features meaningful. With linear models and a noisy
self-reported target, further gains are small, and differences below the
fold standard deviation ($\approx0.005$) should not be over-interpreted. The
main limitations are the linear decision boundary, median imputation, which
ignores relations between questions, and our hand-written metadata, which may
contain coding mistakes. Natural next steps are interaction features between
strong predictors (age, hypertension, diabetes), and models with a non-linear
decision boundary.

\end{document}
